\subsection{Learning algorithm in cEBMF form}
\label{sec:learningalgorithm}
The implementation retains cEBMF's residual-based reduction to normal means,
but separates latent-state transitions from estimation of the conditional
prior parameters. We describe a single matrix \(Y\in\R^{N\times P}\), as in
the tree example. The coupled-modality version uses the same steps with all
loading nodes in the specified directed graph.
The equations below give the unpenalized case. The regularized target and
the explicit ELBO diagnostic are defined in Section~\ref{sec:implementedelbo}.

\paragraph{The ordinary cEBMF iteration.}
With fixed covariates and a factorized variational distribution, let
\(\bar L_{ik}=\E_q L_{ik}\), \(\bar F_{tk}=\E_q F_{tk}\), and
\(\overline{F^2}_{tk}=\E_q F_{tk}^2\). Write
\(w_{it}=M_{it}\tau_{it}\), where \(M_{it}\) indicates observation, and form
\begin{align}
\bar E_{it,-k}&=Y_{it}-\sum_{h\ne k}\bar L_{ih}\bar F_{th},\\
a^{\rm VI}_{ik}&=\sum_t w_{it}\overline{F^2}_{tk},&
b^{\rm VI}_{ik}&=\sum_t w_{it}\bar E_{it,-k}\bar F_{tk}.
\label{eq:algorithmvi}
\end{align}
All sums involving \(Y\) are over observed entries. For positive \(a\),
the resulting normal-means observation is \(\widehat l_{ik}=b_{ik}/a_{ik}\)
with standard error \(s_{ik}=a_{ik}^{-1/2}\). A conditional EBNM fit estimates
the loading prior and returns first and second posterior moments. The factor
update exchanges \(L\) and \(F\) and uses the updated loading moments.
Thus the ordinary loop is: residualize, fit the loading normal-means model,
fit the factor normal-means model, and repeat over factors. Its criterion is
the variational objective when local fitting steps optimize that criterion.

\paragraph{The corresponding joint latent update.}
The sampling state contains realized \(L,F\), mixture labels, and any HMM
paths. Its residual and likelihood statistics use current values:
\begin{align}
E_{it,-k}&=Y_{it}-\sum_{h\ne k}L_{ih}F_{th},\\
a^L_{ik}&=\sum_t w_{it}F_{tk}^2,&
b^L_{ik}&=\sum_t w_{it}E_{it,-k}F_{tk},
\label{eq:algorithmloadingstats}\\
a^F_{tk}&=\sum_i w_{it}L_{ik}^2,&
b^F_{tk}&=\sum_i w_{it}E_{it,-k}L_{ik}.
\label{eq:algorithmsamplingstats}
\end{align}
The residual is refreshed after each column update. These are conditional
likelihood statistics: posterior second moments must not be substituted
for the squared realized values. The natural-parameter form
\(\exp(-av^2/2+bv)\) also handles \(a=b=0\) without dividing by zero.

For a scalar node \(u\), let \(P_{iu}\) contain its current latent parents
and any fixed covariates. At fixed prior parameters, draw a value and label
from the normal-means posterior
\begin{equation}
q_{0,iu}(dv,z)\ \propto\
\exp\{-a_{iu}v^2/2+b_{iu}v\}\,
g_{\theta u}(dv,z\mid P_{iu}).
\label{eq:algorithmproposal}
\end{equation}
This uses the posterior-calculation part of EBNM; it does \emph{not} refit
the prior for each proposal. The full conditional additionally contains
the direct children's augmented prior terms \(D_{iu}(v)\) from
\eqref{eq:implementedmessage}. Accept the proposed \((v',z')\) with
\begin{equation}
\alpha=\min\{1,\exp[D_{iu}(v')-D_{iu}(v)]\}.
\label{eq:algorithmcorrection}
\end{equation}
All other values and labels remain fixed during this move. The parent
proposal density cancels because the full conditional is proportional to
\(q_{0,iu}(dv,z)e^{D_{iu}(v)}\). Without children, the acceptance probability
is one. Fixed side information enters both own and child priors, but is
never updated or assigned an additional likelihood. Independent row moves
are vectorized within a column; columns are updated sequentially.

\paragraph{The empirical-Bayes parameter update.}
After \(S\) complete sweeps at fixed parameters, retain the full states
\(\{L^{(s)},F^{(s)},Z^{(s)},\ldots\}_{s=1}^S\). For each learned loading
prior, reduce
\begin{equation}
J^L_{k,t}(\theta_k)=-\frac{1}{SN}\sum_{s=1}^S\sum_{i=1}^N
\log g_{\theta k}\!\left(L_{ik}^{(s)},Z_{ik}^{(s)}
\mid X_{l,i},L_{i,<k}^{(s)}\right).
\label{eq:algorithmlearningloss}
\end{equation}
Omit \(L_{i,<k}^{(s)}\) on an axis without self-covariates. The factor-axis
objective exchanges \(N,L,X_l\) with \(P,F,X_f\). Parents, responses and
labels come from the \emph{same} draw. No measurement-error deconvolution is
applied to these latent values. With distinct parameter blocks for different
priors, the complete-data parameter objective separates into these node
objectives; likelihood and fixed-prior terms do not depend on the parameters
being fitted. If both axes are learned, the total negative Monte Carlo
\(Q\), averaged over draws, is
\(N\sum_kJ^L_{k,t}+P\sum_kJ^F_{k,t}\), up to a constant.
The per-axis normalizations do not change separate node optima.

\textbf{Relationship to the package interface.}
\par\smallskip
An effective \texttt{self\_row\_cov=True} or
\texttt{self\_col\_cov=True} selects this algorithm.
\texttt{fit(T)} performs the learning rounds and retained sampling.
The loop controls are supplied through \texttt{joint\_kwargs}.
Its keys are \path{sweeps_per_round}, \path{steps}, \path{burnin},
\path{draws} and \path{thin}.
For learned priors, \path{n_epochs} overrides \path{steps} on its axis,
and \path{batch_size} specifies the training mini-batches.
\texttt{iter\_once()} performs one fixed-parameter sweep, rather than the
entire E/M round. The implemented scan updates all loading columns before
all factor columns; the ordinary variational route alternates loading and
factor updates within each rank-one component. In the coupled ATAC--RNA
solver, step 2(a)(i) visits the combined loading DAG, including RNA child
terms in ATAC moves; step 2(a)(ii) updates each modality's feature factors.

\paragraph{What a decreasing curve establishes.}
For the tree notebook, \(J_t=\sum_kJ^L_{k,t}\) is the plotted prior-learning
loss. The retained parameter state satisfies
\(J_{k,t}(\theta_{k,t+1})\leq J_{k,t}(\theta_{k,t})\) on the same saved
draws. Intermediate Adam steps need not decrease it. Draws change between
rounds, so the before/after comparison is justified within a round, not as
an exact trajectory of marginal likelihood across rounds. Fixed-parameter
Markov kernels preserve the posterior; finite chains do not provide exact
E-steps or a guarantee of marginal-likelihood ascent. The implementation
runs a requested number of rounds, not an evidence-based stopping rule.
Afterward, the sampled log joint fluctuates and is not an ELBO. A separate
ELBO diagnostic, including entropy, is now evaluated as described in
Section~\ref{sec:implementedelbo}. Held-out
prediction error is a separate evaluation metric. These distinctions
preserve the familiar cEBMF computational structure without attributing
variational coordinate-ascent guarantees to the joint sampler.

\clearpage
\noindent\textbf{Algorithm 1: residual-based joint cEBMF with partial Monte Carlo EM}
\par\smallskip\hrule\smallskip
\textbf{Input.} Observed matrix, mask and any known noise variances; fixed
row/column covariates; prior
families and factor order; learning rounds \(T\), sweeps per round \(S\),
training epochs \(B\), optional batch size, burn-in \(b_0\), retained draws
\(D\), and thinning \(h\).
\begin{enumerate}[leftmargin=*,label=\textbf{\arabic*.},itemsep=5pt]
\item \textbf{Initialize.} Obtain \(L,F\) from a preliminary independent
fit, or use supplied initial factors. Fix the resulting rank, order,
observation precisions, preprocessing and non-neural prior parameters.
Initialize and pretrain the conditional-prior networks and establish valid
component labels and HMM paths. Pretraining on initial estimates is an
initialization device, not additional data in the joint likelihood.
\item \textbf{For learning rounds \(t=1,\ldots,T\):}
  \begin{enumerate}[label=(\alph*),leftmargin=*,itemsep=3pt]
  \item \textbf{Approximate E-step.} Hold all parameters fixed. Repeat \(S\)
  times, continuing from the current chain state:
    \begin{enumerate}[label=(\roman*),leftmargin=*,itemsep=2pt]
    \item For \(k=1,\ldots,K\), residualize and update \(L_{:k}\) using
    \eqref{eq:algorithmloadingstats} and the appropriate prior kernel.
    Neural scalar priors use \eqref{eq:algorithmproposal}--\eqref{eq:algorithmcorrection}.
    \item For \(k=1,\ldots,K\), update \(F_{:k}\) analogously using the
    new loading state. A fixed Gaussian-mixture prior uses direct conditional
    draws. An HMM prior uses a whole-path forward-filtering/backward-sampling
    update and conditional value draws. Use the same kernel choices on
    either axis; an HMM ignores covariates on its own axis.
    \item Save the full joint state after the complete sweep.
    \end{enumerate}
  \item \textbf{Generalized M-step.} Hold these \(S\) states fixed.
  For each neural prior on each axis, run up to \(B\) epochs of Adam on
  \eqref{eq:algorithmlearningloss} using full or mini-batches. After each
  epoch evaluate the full saved-draw loss; retain the best finite parameter
  state, including the starting state. Record the loss before and after.
  Continue the chain from its last realized state for the next round.
  \end{enumerate}
\item \textbf{Posterior sampling.} Freeze all fitted parameters. Discard
\(b_0\) additional sweeps, then retain \(D\) full states, one every \(h\)
sweeps. No prior learning, noise update or factor pruning occurs here.
\item \textbf{Return.} Report marginal moments and the signal estimate
\[
\widehat X=\frac1D\sum_{d=1}^D L^{(d)}(F^{(d)})^\top,
\]
together with the prior-learning losses and sampling diagnostics.
\end{enumerate}
\smallskip\hrule\medskip

\clearpage
